\documentclass[]{article}

%opening
\title{}
\author{}
\usepackage{minted}

\begin{document}

\maketitle


\label{cha:1}
textit{I will intruduce concepts like foas, hoas, phoas and verification condition generation. Then I will motivate using phoas for verification condition generation}

\section{Representing abstract syntax}

When a program in a programming language is transforming the syntax of some other language, it often uses an abstract representation of the syntax in that programming language. Examples of this are compilers, interpreters, typecheckersand code generators. Since many languages have variables, the program transforming the syntax must have a way to deal with variable binding and scoping. Often this involves implementing substitution and finding ways to handle shadowing. This often can lead to bugs, for example in \cite{captureAvoiding} a study was done on 9 DSL implementations and 8 of them were prone to variable capture. Some abstract representations make these variable bindings easier to deal with. I will discuss different representations of abstract syntax and discuss their advantages and disadvantages, focussing on different ways to handle variable bindings. In what follows I will call the language we're representing the syntax of the object language and the language we're representing the object language in the meta-language. As the examples, I will represent the object language of the untyped lambda calculus in the meta-language Haskell.

\subsection{First Order Abstract Syntax (Foas)}
In a first order abstract syntax, variables are represented with names. 



\subsubsection{Nominal names}
The first first order approach is one that uses names nominally, I will discuss an approach that uses a \mintinline{haskell}{String} for these names.

\begin{figure}
	\centering
	\begin{minted}
		data Lam = Var String
		| Lambda String Lam
		| App Lam Lam
		
		apply :: Lam
		apply = Lambda “f” (Lambda “x” (App (Var “f”) (Var “x”)))
		
		foasI :: Lam
		foasI = Lambda "x" (Var "x")
	\end{minted}
	\caption{Foas representation of the untyped lambda calculus}
	\label{fig:foas_lam}
\end{figure}

The examples represent the identity function $(\lambda x. x)$ and function application $(\lambda f. \lambda x. f  x ) $

When a variable is introduced with a lambda, a name for it is given as a \mintinline{haskell}{String}. When a variable is used, this string is used to reference which introduction corresponds to which variable.

Advantages of this approach are that each variable has a unique and explicit identity. Additionally the name of a variable is the same everywhere it occurs. The syntax can easily be printed out and used for intentional analysis \cite{atkey2009unembedding}. To represent syntax in this way, the meta-language does not require any advanced features such as higher order functions. 

Some disadvantages of this approach are that it is possible to represent ill-formed syntax with this approach. For example, \mintinline{haskell}{Lamba "x" (Var "y")} can be respresented in this Foas approach, but the variable y was never introduced. Additionally there are more representations possible for alpha equivalent terms. For example, \mintinline{haskell}{Lamba "x" (Var "x")} and \mintinline{haskell}{Lamba "y" (Var "y")} represent the same function but have different representations.


\subsubsection{De Bruijn Indices}
The previous approach had the disadvantage that alpha-equivalent terms can have multiple representations. De Bruijn indecis try to solve this by giving no choice for what the name is when a variable is introduced. When a variable is used it is the natural number that measures the amount of lambda abstractions done between the introduction of the variable and the usage.

\begin{figure}
	\centering
	\begin{minted}{haskell}
		data DBTerm = DBVar Int
		| DBLam DBTerm
		| DBApp DBTerm DBTerm
		deriving (Show,Eq)
		
		
		deBruinI :: DBTerm
		deBruinI = DBLam (DBVar 0)
		
		apply :: DBTerm
		apply = DBLam (DBLam (DBApp (DBVar 1) (DBVar 0)))
		
	\end{minted}
	\caption{De Bruijn Index representation of the untyped lambda calculus}
	\label{fig:DeBruijn_Lam}
\end{figure}


Advantages of this approach are that all alpha equivalent terms have a unique representation. Additionally each variable has a explicit identity. Just like the nominal names foas approach, the syntax can easily be printed out, used for intentional analysis \cite{atkey2009unembedding}and the meta-language does not need to support higher order functions. 


A disadvantage of this approach is that you can the same variable might have different names throughout depending on where it occurs. For example the term $(\lambda x. x\  (\lambda f. x))$ would be represented as \mintinline{haskell}{DBLam (DBApp (Var 0) (DBLam (Var 1)))} where x is represented as both \mintinline{haskell}{Var 0 } and \mintinline{haskell}{Var 1}. Another disadvantage would be that this representation can still create ill-formed terms. For example the term \mintinline{haskell}{DBLam (DBVar 7)} references a term that isn't introduced. 



\subsection{Higher order abstract syntax (HOAS)}

Higher Order Abstract Syntax was introduced by \cite{HOAS} of representing abstract syntax in a meta-language supporting higher order functions. It uses the binding construct of the meta-language to represent the binding.

\begin{figure}
	\centering
	\begin{minted}
		data Lam = App Lam Lam
		| Lambda (Lam -> Lam)
		
		hoasI :: Lam
		hoasI = Lambda (\x->x)
		
		apply :: Lam
		apply = Lambda (\f -> Lambda (\x -> App f x))
	\end{minted}
	\caption{Hoas representation of the untyped lambda calculus}
	\label{fig:hoas_lambda}
\end{figure}

The advantage of using this approach is that there are no variables that can occur out of scope. According to https://jesper.sikanda.be/posts/1001-syntax-representations.html substitution can be implemented simply by using function application on the meta-level and alpha equivalent terms are equal.

Even though ill-formed terms are no longer representable in this dataype, there are still members of this datatype that do not represent terms of the lambda calculus called "exotic types" Consider the following example from \cite{atkey2009unembedding}
\begin{figure}
	\centering
	\begin{minted}
		
		exotic :: Lam
		exotic = Lambda (\x -> case x of 
		| Lam f -> x
		| App a b -> Lam (\x -> x))
	\end{minted}
	\caption{Example of an exotic term in Hoas}
	\label{fig:exotic}
\end{figure}
Since Haskell can pattern match on the inductive datatype of Lam, we create a function that is not possible in the lambda calculus, where pattern matching is not possible. Additionaly this syntax representation is inconventient to transform and do functional programming with. For example it is not clear how to count the amount of variable occurences. Another example that is inconvenient is that it is not clear how to convert a Hoas syntax representation to a Foas one. 






\subsection{Parametric Higher Order Abstract Syntax}

(information from eg. Parametric Higher-Order Abstract Syntax for Mechanized
Semantics)

pro: convertable to FOAS, functional programming, geen exotic terms


con: phoas, inspecting bodies

\subsection{Translations between representations}
Describe translations phoas<->foas and phoas-> hoas such as described in Unembedding Domain-Specific Languages.


\section{monads?}
(I'm not sure if this is necessary but it is something most other students would not have as background knowledge)
Describe option, reader, state and continuation monads.
\denis{yes fine, but be brief. Not 3 pages. Use it as a step up to explain specification monads}

\section{Verification condition generation}
Describe what a verification condition is and how it can be computed by recursively calculating the weakest precondition.

\subsection{Wstore and its combinators}
Describe the Wstore monad and assert, demonic and angelic choice combinators.
\subsection{generating verification conditions}
Describe how to use these combinators to generate the weakest precondition.



(Verified Symbolic Execution with Kripke Specification
Monads (and No Meta-programming)




\end{document}
